\begin{exercisechunk}
\item \textbf{What signed correlation controls.}
\label[exercise]{ex:l7-correlation-loss}
\exerciseprofile{2}{2}{2}{\exproof\quad\excalculation\quad\exconstruction\quad\exinterpretation}{%
Relate signed correlation to classification and to the gradients of standard
surrogate losses.}
Under the fixed-parity law $P_S$, recall
\[
\cC(f)=-\E_{P_S}[Yf(X)].
\]
\begin{enumerate}[(a),beginpenalty=10000]
\item Let $h:\{-1,1\}^d\to\{-1,1\}$. Prove
\[
\cC(h)=2R_{01}(h)-1,
\qquad
R_{01}(h)=P_S(h(X)\ne Y).
\]
\item For $\varepsilon>0$, set $f_\varepsilon=\varepsilon\chi_S$. Compute
$\cC(f_\varepsilon)$ and $R_{01}(\operatorname{sign}(f_\varepsilon))$.
Explain why a correlation value close to zero does not determine the
classification risk of the sign of a real-valued score.
\item For any score $f:\{-1,1\}^d\to\R$, derive
\[
\E_{P_S}[(f(X)-Y)^2]
=1+\E_{P_S}[f(X)^2]+2\cC(f).
\]
\item Derive
\[
\E_{P_S}[\log(1+e^{-Yf(X)})]
=\E_{P_S}\!\left[\log\!\left(2\cosh\frac{f(X)}2\right)\right]
+\frac12\cC(f).
\]
If $f=f_\theta$ is differentiable in $\theta$, differentiate the identities
in parts (c) and (d) and identify every term in addition to
$\nabla_\theta\cC(f_\theta)$.
\end{enumerate}
\begin{hints}
For (a), determine the value of $Yh(X)$ on the correct and incorrect events.
For (d), verify the pointwise identity
\[
\log(1+e^{-z})=\log\!\left(2\cosh\frac z2\right)-\frac z2.
\]
\end{hints}
\end{exercisechunk}
